\begin{thebibliography}{29}
\bibitem{ref1} Miklós Ajtai and Avi Wigderson. Deterministic simulation of probabilistic constant depth circuits. Adv. Comput. Res., 5:199–222, 1989.
\bibitem{ref2} Paul Beame, Russell Impagliazzo, and Srikanth Srinivasan. Approximating ac\textasciicircum{}0 by small height decision trees and a deterministic algorithm for \#ac\textasciicircum{}0sat. In Computational Complexity Conference, pages 117–125. IEEE Computer Society, 2012.
\bibitem{ref3} Mark Braverman. Polylogarithmic independence fools AC 0 circuits. J. ACM, 57(5):28:1–28:10, 2010. doi:10.1145/1754399.1754401.
\bibitem{ref4} Ruiwen Chen, Rahul Santhanam, and Srikanth Srinivasan. Average-case lower bounds and satisfiability algorithms for small threshold circuits. Theory Comput., 14(1):1–55, 2018. doi:10.4086/toc.2018.v014a009.
\bibitem{ref5} Anindya De, Omid Etesami, Luca Trevisan, and Madhur Tulsiani. Improved pseudoran- dom generators for depth 2 circuits. In Maria J. Serna, Ronen Shaltiel, Klaus Jansen, and José D. P. Rolim, editors, Approximation, Randomization, and Combinatorial Optimiza- tion. Algorithms and Techniques, 13th International Workshop, APPROX 2010, and 14th International Workshop, RANDOM 2010, Barcelona, Spain, September 1-3, 2010. Proceed- ings, volume 6302 of Lecture Notes in Computer Science, pages 504–517. Springer, 2010. doi:10.1007/978-3-642-15369-3\_38.
\bibitem{ref6} Dean Doron, Raghu Meka, Omer Reingold, Avishay Tal, and Salil P. Vadhan. Monotone branching programs: Pseudorandomness and circuit complexity. Electron. Colloquium Comput. Complex., page 18, 2021. URL: https://eccc.weizmann.ac.il/report/2021/018.
\bibitem{ref7} Michael A. Forbes and Zander Kelley. Pseudorandom generators for read-once branching pro- grams, in any order. In Mikkel Thorup, editor, 59th IEEE Annual Symposium on Foundations of Computer Science, FOCS 2018, Paris, France, October 7-9, 2018, pages 946–955. IEEE Computer Society, 2018. doi:10.1109/FOCS.2018.00093.
\bibitem{ref8} Oded Goldreich and Avi Wigderson. On the size of depth-three boolean circuits for computing multilinear functions. Electron. Colloquium Comput. Complex., page 43, 2013. URL: https: //eccc.weizmann.ac.il/report/2013/043.
\bibitem{ref9} Parikshit Gopalan, Raghu Meka, and Omer Reingold. DNF sparsification and a faster deterministic counting algorithm. Comput. Complex., 22(2):275–310, 2013. doi:10.1007/ s00037-013-0068-6.
\bibitem{ref10} Parikshit Gopalan, Raghu Meka, Omer Reingold, Luca Trevisan, and Salil P. Vadhan. Better pseudorandom generators from milder pseudorandom restrictions. In 53rd Annual IEEE Symposium on Foundations of Computer Science, FOCS 2012, New Brunswick, NJ, USA, October 20-23, 2012, pages 120–129. IEEE Computer Society, 2012. doi:10.1109/FOCS.2012. 77.
\bibitem{ref11} Elad Haramaty, Chin Ho Lee, and Emanuele Viola. Bounded independence plus noise fools products. SIAM J. Comput., 47(2):493–523, 2018. doi:10.1137/17M1129088.
\bibitem{ref12} Prahladh Harsha and Srikanth Srinivasan. On polynomial approximations to AC. Random Struct. Algorithms, 54(2):289–303, 2019. doi:10.1002/rsa.20786.
\bibitem{ref13} Johan Håstad. On the correlation of parity and small-depth circuits. SIAM J. Comput., 43(5):1699–1708, 2014. doi:10.1137/120897432.
\bibitem{ref14} John Hastad. Almost optimal lower bounds for small depth circuits. Adv. Comput. Res., 5:143–170, 1989.
\bibitem{ref15} Pooya Hatami, William Hoza, Avishay Tal, and Roei Tell. Fooling constant-depth threshold circuits. Electron. Colloquium Comput. Complex., page 2, 2021. URL: https://eccc.weizmann. ac.il/report/2021/002.
\bibitem{ref16} Russell Impagliazzo, William Matthews, and Ramamohan Paturi. A satisfiability algorithm for ac0 . In SODA, pages 961–972. SIAM, 2012.
\bibitem{ref17} Russell Impagliazzo, Raghu Meka, and David Zuckerman. Pseudorandomness from shrinkage. In 53rd Annual IEEE Symposium on Foundations of Computer Science, FOCS 2012, New Brunswick, NJ, USA, October 20-23, 2012, pages 111–119. IEEE Computer Society, 2012. doi:10.1109/FOCS.2012.78.
\bibitem{ref18} Zander Kelley. An improved derandomization of the switching lemma. In Samir Khuller and Virginia Vassilevska Williams, editors, STOC ’21: 53rd Annual ACM SIGACT Symposium on Theory of Computing, Virtual Event, Italy, June 21-25, 2021, pages 272–282. ACM, 2021. doi:10.1145/3406325.3451054.
\bibitem{ref19} Chin Ho Lee and Emanuele Viola. More on bounded independence plus noise: Pseudorandom generators for read-once polynomials. Theory Comput., 16:1–50, 2020. doi:10.4086/toc. 2020.v016a007.
\bibitem{ref20} Nathan Linial, Yishay Mansour, and Noam Nisan. Constant depth circuits, fourier transform, and learnability. J. ACM, 40(3):607–620, 1993.
\bibitem{ref21} Michael Luby and Boban Velickovic. On deterministic approximation of DNF. Algorithmica, 16(4/5):415–433, 1996. doi:10.1007/BF01940873.
\bibitem{ref22} Raghu Meka, Omer Reingold, and Avishay Tal. Pseudorandom generators for width-3 branching programs. In Moses Charikar and Edith Cohen, editors, Proceedings of the 51st Annual ACM SIGACT Symposium on Theory of Computing, STOC 2019, Phoenix, AZ, USA, June 23-26, 2019, pages 626–637. ACM, 2019. doi:10.1145/3313276.3316319.
\bibitem{ref23} Raghu Meka and David Zuckerman. Pseudorandom generators for polynomial threshold functions. SIAM J. Comput., 42(3):1275–1301, 2013. doi:10.1137/100811623.
\bibitem{ref24} Noam Nisan and Avi Wigderson. Hardness vs randomness. J. Comput. Syst. Sci., 49(2):149–167, 1994. doi:10.1016/S0022-0000(05)80043-1.
\bibitem{ref25} Rocco A. Servedio and Li-Yang Tan. Improved pseudorandom generators from pseudorandom multi-switching lemmas. In Dimitris Achlioptas and László A. Végh, editors, Approximation, Randomization, and Combinatorial Optimization. Algorithms and Techniques, APPROX/RAN- DOM 2019, September 20-22, 2019, Massachusetts Institute of Technology, Cambridge, MA, USA, volume 145 of LIPIcs, pages 45:1–45:23. Schloss Dagstuhl - Leibniz-Zentrum für Inform- atik, 2019. doi:10.4230/LIPIcs.APPROX-RANDOM.2019.45.
\bibitem{ref26} Avishay Tal. Tight bounds on the fourier spectrum of AC0. In Ryan O’Donnell, editor, 32nd Computational Complexity Conference, CCC 2017, July 6-9, 2017, Riga, Latvia, volume 79 of LIPIcs, pages 15:1–15:31. Schloss Dagstuhl - Leibniz-Zentrum für Informatik, 2017. doi: 10.4230/LIPIcs.CCC.2017.15.
\bibitem{ref27} Luca Trevisan and Tongke Xue. A derandomized switching lemma and an improved derandom- ization of AC0. In Proceedings of the 28th Conference on Computational Complexity, CCC 2013, K.lo Alto, California, USA, 5-7 June, 2013, pages 242–247. IEEE Computer Society, 2013. doi:10.1109/CCC.2013.32.
\bibitem{ref28} Salil P. Vadhan. Pseudorandomness. Found. Trends Theor. Comput. Sci., 7(1-3):1–336, 2012. doi:10.1561/0400000010.
\bibitem{ref29} Emanuele Viola. AC0 unpredictability. ACM Trans. Comput. Theory, 13(1):5:1–5:8, 2021. doi:10.1145/3442362.
\end{thebibliography}
